\section{Invariant Subspaces}

\subsection{Eigenvalues}

% 5.1
\begin{mydef} [operator]
  A \lm from a \vs to itself is called an \qt{operator}.
\end{mydef}

(Suppose $T\in \linmap(V)$; then $\left.T\right|_{V_{k}}$ may not be an operator on a subspace $V_k$.)

% 5.2
\begin{mydef} [invariant subspace]
  Let $T\in \linmap(V)$. A subspace $U$ of $V$ is called \qt{invariant under $T$} if $Tu \in U$ for all $u \in U$. \\
  Thus, $U$ is invariant under $T$ if $\left.T\right|_{U}$ is an operator on $U$.
\end{mydef}

\begin{example}
  Let $T\in \linmap(\mathcal{P}(\mathbb{R}))$ such that $Tp=p'$. Let $U=\mathcal{P}_4(\mathbb{R}) \subseteq \mathcal{P}(\mathbb{R})$. Then $U$ is invariant under $T$
  because if $p \in U$, $\deg p \leq 4$ and $\deg (p') \leq 3$.
\end{example}

\begin{example}
  Let $T\in \linmap(V)$. Then $\{0\}, V, \mynull T, \myrange T$ are all invariant. \\
  (If $T$ is invertible, then $\mynull T = \{0\}$ and $\myrange T=V$.)
\end{example}

\bfemph{Invariant subspaces of dimension one: }

Take any $v\in V, v\neq 0$, and let
\begin{equation}
  U :\equiv \{  \lambda v \mid \lambda \in \myF \} = \myspan{v},
\end{equation}
then $U$ is a one\-/dimensional subspace of $V$.

If $U$ is invariant under an operator $T \in \linmap (V)$, then $Tv \in U \implies \exists \lambda \in \myF: Tv = \lambda v$. \\
Conversely, if $Tv = \lambda v$, $\lambda \in \myF$, then $\myspan{v}$ is a one-dimensional subspace of $V$ invariant under $T$.

% 5.5
\begin{mydef} [eigenvalue]
  Suppose $T\in \linmap (V)$. A scalar $\lambda \in \myF$ is called an \qt{eigenvalue of $T$} if there exists a nonzero (i.e., $v \neq 0$) $v \in V$ such that
  \begin{equation}
    Tv = \lambda v.
  \end{equation}
\end{mydef}

% 5.7
\setcounter{thm}{6}
\begin{thm} [equivalent conditions to be an eigenvalue]
  \label{thm: equivalent conditions to be an eigenvalue}
  Suppose $V$ is finite-dimensional. Then the following are equivalent for $T \in \linmap(V)$ and $\lambda \in \myF$:
  \begin{enumerate}[label=\textbf{(\alph*)}]
    \item $\lambda$ is an eigenvalue of $T$. \label{first}
    \item $T-\lambda I$ is not injective. \label{second}
    \item $T-\lambda I$ is not surjective. \label{third}
    \item $T-\lambda I$ is not invertible. \label{forth}
  \end{enumerate}
\end{thm}
\begin{prf}
  Conditions \ref{first} and \ref{second} are equivalent because the eigenvector $v$ is a solution to $Tv=\lambda v$, which is equivalent to $(T-\lambda I)v=0$. Hence, $\dim \mynull (T-\lambda I) > 0$, which implies that $T-\lambda I$ is not injective.
  \ref{second}, \ref{third}, and \ref{forth} are therefore equivalent by
  \ref{thm: injectivity is equivalent to surjectivity}.
\end{prf}

\setcounter{thm}{7}
\begin{mydef} [eigenvector]
  Let $T\in \linmap(V)$. A vector $v \in V$ is called an \qt{eigenvector} of $T$ corresponding to $\lambda$ if $v\neq 0$ and $Tv = \lambda v$.
  In other words:
  A vector $v\in V, v \neq 0$ is an eigenvector corresponding to $\lambda \iff v \in \mynull(T-\lambda I_V)$.
\end{mydef}

% 5.12
\setcounter{thm}{10}
\begin{thm}[linearly independent eigenvectors]
  \label{thm: linearly independent eigenvectors}
  Every list of eigenvectors of $T$ corresponding to distinct eigenvalues of $T$ is linearly independent.
\end{thm}
\begin{prf}
  Suppose the desired result is false. Then there exists a smallest list of length $m$ of linearly dependent eigenvectors $v_1, \dots, v_m$ with distinct eigenvalues $\lambda_1, \dots, \lambda_m$. Since an eigenvector is unequal to the zero vector, $m$ must be $\geq 2$.

  Because of the minimality of $m$ and because our list is linearly dependent:
  \begin{equation}
    \exists a_1, \dots, a_m\text{, where not all $a$'s are equal to }0\text{, such that } a_1 v_1 + \cdots + a_m v_m = 0.
  \end{equation}

  Now, we apply $(T-\lambda_m I)$ on both sides of the equation and get
  \begin{equation}
    \begin{gathered}
      (a_1 \lambda_1 v_1 - a_1 \lambda_m v_1)
      + \cdots + \\
      (a_{m-1} \lambda_{m-1} v_{m-1} - a_{m-1} \lambda_{m} v_{m-1}) +
      (\underbrace{a_m \lambda_m v_m -a_m \lambda_m v_m}_{=0}) =0,
    \end{gathered}
  \end{equation}

  which is the same as:
  \begin{equation}
    a_1 \underbrace{(\lambda_1 - \lambda_m)}_{\neq 0} v_1 + \cdots + a_{m-1} \underbrace{(\lambda_{m-1}-\lambda_{m})}_{\neq 0} v_{m-1}=0,
  \end{equation}

  By the minimality of $m$, the shorter list $v_1, \ddd, v_{m-1}$ is linearly independent, so the equation above forces $a_k(\lambda_k - \lambda_m) = 0$ for each $k \in \{1, \ddd, m-1\}$. Because the eigenvalues are distinct, $\lambda_k - \lambda_m \neq 0$, and hence $a_1 = \cdots = a_{m-1} = 0$. The original equation then reduces to $a_m v_m = 0$, and since $v_m \neq 0$ we also get $a_m = 0$. This contradicts the choice of the $a$'s as not all zero. Therefore, no such linearly dependent list of eigenvectors can exist.
\end{prf}

\begin{thm}[number of eigenvalues]
  \label{thm: operator cannot have more eigenvalues than dimension of vector space}
  Each operator $T \in \linmap(V)$ on $V$ has at most $\dim V$ distinct eigenvalues.
\end{thm}
\begin{prf}
  Let $T \in \linmap(V)$. Suppose $\lambda_1, \ddd, \lambda_m$ are distinct eigenvalues of $T$ with corresponding eigenvectors $v_1, \ddd, v_m \in V$. Hence, \ref{thm: linearly independent eigenvectors} implies that the list $v_1, \ddd, v_m$ is linearly independent. Thus, \ref{thm: length of linearly independent list at most length of spanning list}
  implies that $m \leq \dim V$.
\end{prf}


\subsection{Polynomials Applied to Operators}

% 5.13
\setcounter{thm}{12}
\begin{mydef} [notation $T^m$]
  Let $T \in \linmap(V)$ and $m\in \nat^{+}$:
  \begin{itemize}
    \item $T^{m} :\equiv \underbrace{T \cdots T}_{\text{$m$ times}}$ or $T^{m} :\equiv T^{m-1} \cdot T$ such that $T^{m} \in \linmap(V)$;
    \item $T^0 :\equiv I_V$;
    \item If $T$ is invertible with inverse $T^{-1}$, then $T^{-m}\in \linmap(V)$ is defined by $T^{-m} :\equiv (T^{-1})^m$.
  \end{itemize}
\end{mydef}
$T^m T^n = T^{m+n}$ and $(T^m)^n=T^{mn}$ hold for $m,n \in \mathbb{Z}$ if $T$ is invertible, and for $m,n \in \mathbb{N}$ if $T$ is not invertible.

\begin{mydef} [notation $p(T)$]
  For $\mathcal{P} (\myF) \ni p(z) = a_0+a_1z+a_2z^2+\cdots+a_mz^m$
 and
  $T \in \linmap (V)$, we define:
  \begin{equation}
    p(T) :\equiv a_0 I + a_1 T + a_2 T^2 + \cdots + a_m T^m.
  \end{equation}

  Note that $p(T) \in \linmap(V)$.
\end{mydef}

%TODO: example 5.15??

%TODO: example 5.16??

\setcounter{thm}{16}
\begin{thm} [multiplicative properties]
  \label{multiplicative-properties}
  Suppose $p,q \in \mathcal{P} (\myF)$ and $T\in \linmap (V)$. Then
  \begin{equation}
    (p q)(T) = p(T) q(T) = q(T)p(T).
  \end{equation}
\end{thm}
\begin{prf} When a product of polynomials is expanded using the distributive property, it does not matter if the symbol is $z$ or $T$. Suppose $p(z) = \sum_{j=0}^{m} a_j z^j$ and $q(z)=\sum_{k=0}^{n} b_k z^k$ for all $z \in \myF$. Then % wether?
  \begin{equation}
    (pq)(z) = \sum_{j=0}^{m} \sum_{k=0}^{n} a_j b_k z^{j+k} \myand
    (pq)(T) = \sum_{j=0}^{m} \sum_{k=0}^{n} a_j b_k T^{j+k}
    = \sum_{j=0}^{m} a_j T^j \sum_{k=0}^{n}  b_k T^k,
  \end{equation}
  which is the same as $(pq)(T) = p(T)q(T)$ as well as $(pq)(T) = q(T)p(T)$.
\end{prf}

\begin{thm}[null space and range of $p(T)$ are invariant under $T$]
  \label{thm: null space and range of p(T) are invariant under T}
  For every $T \in \linmap(V)$ and every $p \in \mathcal{P} (\myF)$,
  \begin{equation}
    \mynull p(T) \text{ and } \myrange p(T) \text{ are invariant under } T.
  \end{equation}
\end{thm}
\begin{prf}
  \StepOne First, let $u \in \mynull p(T)$. Then,
  \begin{equation}
    p(T)u = 0.
  \end{equation}

  Since $p(T)$ commutes with $T$ --- which is \ref{multiplicative-properties} applied to the polynomials $p(z)$ and $z$, whose product is the same in either order ---
  \begin{equation}
    \begin{aligned}
      (p(T))(Tu) &= (p(T)T)(u) \\
      &= (Tp(T))(u) \\
      &= T(p(T)u) \\
      &= T(0) = 0.
    \end{aligned}
  \end{equation}

  Hence, $Tu \in \mynull p(T)$.

  \StepTwo Second, let $u \in \myrange p(T)$. Then, $u = p(T)v$ for some $v \in V$. Thus,
  \begin{equation}
    \begin{aligned}
      Tu = T(p(T)v) = p(T)(Tv),
    \end{aligned}
  \end{equation}
  so $Tu \in \myrange p(T)$.
\end{prf}

\begin{ainote}[Commuting is the whole content of this theorem]
  Strip away the notation and both steps use one fact: $T$ commutes with $p(T)$. That is what lets $T$ be moved across $p(T)$ in Step 1 and back again in Step 2.

  So the theorem is really an instance of something more general and worth remembering: \emph{if $S$ and $T$ commute, then $\mynull S$ and $\myrange S$ are invariant under $T$}. The proof is the same three lines, with $p(T)$ replaced by $S$. Polynomials in $T$ are simply the most common source of operators that commute with $T$.

  This is used constantly later on. Every argument that restricts $T$ to $\mynull (T-\lambda I)^k$ or to $\myrange (T - \lambda I)^k$ --- the generalised eigenspaces --- is quoting this result, since those spaces are exactly the null spaces and ranges of polynomials in $T$.
\end{ainote}
